% problem-000347 1 锰锌矿同槽浸出
\section*{1. 锰锌矿同槽浸出}
\textit{下列为分问。}

\subsection*{1-1}
锰结核矿和锌精矿单独酸浸结果很不理想。请通过热⼒学计算，说明锰结核矿和锌精矿同糟酸浸时发⽣化学反应的可⾏性。

\subsection*{1-2}
模拟实验发现，⼆氧化锰和硫化锌同槽酸浸时反应速率很慢，若在酸溶液中加⼊少量的铁屑则能明显使反应速率加快。写出铁进⼈溶液后分别与⼆氧化锰和硫化锌发⽣化学反应的离⼦⽅程式，并简述反应速率加快的原因。

\subsection*{1-3}
研究发现，2种矿物同槽酸浸4⼩时后，锰和锌的浸出率只有\~80%，为了提⾼浸出率，在实际⼯艺中，须将过滤后的滤渣⽤四氯⼄烷处理后再做⼆次酸浸，请简要说明四氯⼄烷的作⽤。

\subsection*{1-4}
锌精矿中常有部分铅共⽣，同槽酸浸⼯艺回收铅的⽅法之⼀是在浸渣中加⼈⾜量的⻝盐⽔使铅溶出，写出溶出铅的化学⽅程式。
